% !TEX root = ../presentation.tex
% =============================================================
%  Subsection — Power transmission (principle definition)
%  Part of ac-transmission; \input from ac-transmission.tex,
%  directly before natural-power.tex.
%
%  Power exchange at the PCC, line parameters, load cases and the
%  load cases.  Figures in ./fig (image4 ... image16).
%  The gamma / Z_w and line-equation frames that used to close this
%  file now live in wave-representation.tex.
% =============================================================
\providecommand{\acfig}[1]{ac-transmission/fig/#1}

\intermezzo{Power transmission}{Principle definition}

\begin{frame}{Power exchange in the AC grid}
  \begin{columns}[c,onlytextwidth]
    \begin{column}{0.44\textwidth}\centering
      \topoimg[width=\linewidth,height=0.5\textheight,keepaspectratio]{\acfig{image4.pdf}}\\[0.3em]
      
    \end{column}
    \begin{column}{0.52\textwidth}\centering
      {\small
      \begin{tabularx}{\linewidth}{@{}l@{\hspace{0.6em}}X@{}}
        $P_S$ & is the export active power of the VSC converter station at PCC\\
        $Q_S$ & is the export reactive power of the VSC converter station at PCC\\
        $U_S$ & is the RMS value of the fundamental frequency component of the line-to-line AC voltage of the VSC converter station at PCC\\
        $U_E$ & is the equivalent AC voltage of the AC grid, corresponding to a SCL at PCC\\
        $R_S$ & is the equivalent resistance of the SCL impedance at PCC\\
        $X_S$ & is the equivalent reactance of the SCL impedance at PCC\\
        $\delta$ & is the phase angle difference between AC voltages $U_S$ and $U_E$\\
      \end{tabularx}}\\[0.3em]
      {\color{atGray}\footnotesize IEC TR 63363-1:2022}
    \end{column}
  \end{columns}
\end{frame}

\begin{frame}{Power exchange in the AC grid}
  \centering
  \topoimg[width=\linewidth,height=0.4\textheight,keepaspectratio]{\acfig{image4.pdf}}\\[0.3em]
  \vspace{1em}     
  $P_S = \frac{U_S}{R_{S}^{2}+X_{S}^{2}} \cdot \left[ \left ( U_S - U_E \cdot \cos \delta \right )  \cdot R_S + U_E \cdot X_S \cdot \sin \delta \right]$\\
  \vspace{1em}
  $Q_S = \frac{U_S}{R_{S}^{2}+X_{S}^{2}} \cdot \left[ \left ( U_S - U_E \cdot \cos \delta \right )  \cdot X_S - U_E \cdot R_S \cdot \sin \delta \right]$\\

\end{frame}

\begin{frame}{Power exchange}
  \centering
  \renewcommand{\arraystretch}{1.25}
  \begin{tabular}{@{}l c c c c@{}}
    \textbf{Parameter} & \textbf{Case 1} & \textbf{Case 2} & \textbf{Case 3} & \textbf{Case 4} \\
    \hline
    $P_S$ & 1000\,MW & 1500\,MW & 1000\,MW & $-1000$\,MW \\
    $U_E$ & 230\,kV & 230\,kV & 230\,kV & 230\,kV \\
    $U_S$ & 232\,kV & 232\,kV & 240\,kV & 232\,kV \\
    $R_S$ & 0.144\,$\Omega$ & 0.144\,$\Omega$ & 0.144\,$\Omega$ & 0.144\,$\Omega$ \\
    $X_S$ & 1.15\,$\Omega$ & 1.15\,$\Omega$ & 1.15\,$\Omega$ & 1.15\,$\Omega$ \\
    $\delta$ & 1.19$^\circ$ & 1.816$^\circ$ & 1.081$^\circ$ & $-1.319^\circ$ \\
    $Q_S$ & 288\,MVAr & 239\,MVAr & 905\,MVAr & 541\,MVAr \\
    \hline
  \end{tabular}

  \vspace{0.6em}
  \begin{itemize}
    \setlength{\itemsep}{0.2em}
    \item Active power transmission requires reactive power.
    \item Higher AC voltage increases the reactive-power demand.
    \item Changing the power-flow direction also changes the reactive-power demand.
  \end{itemize}
\end{frame}

\begin{frame}{Power exchange --- active vs. reactive}
  \begin{columns}[c,onlytextwidth]
    \begin{column}{0.48\textwidth}\centering
      \topoimg[width=\linewidth,height=0.56\textheight,keepaspectratio]{\acfig{image7.png}}\\[0.3em]
      {\footnotesize Active power vs.\ voltage difference}
    \end{column}
    \begin{column}{0.48\textwidth}\centering
      \topoimg[width=\linewidth,height=0.56\textheight,keepaspectratio]{\acfig{image8.png}}\\[0.3em]
      {\footnotesize Active power vs.\ phase difference}
    \end{column}
  \end{columns}
  \vspace{0.3em}
  {\color{atGray}\footnotesize Active power is transmitted via phase difference; reactive power via voltage amplitude.}
\end{frame}

\begin{frame}{Power transmission in the AC grid}
  \begin{columns}[c,onlytextwidth]
    \begin{column}{0.56\textwidth}\centering
      \topoimg[width=\linewidth,height=0.66\textheight,keepaspectratio]{\acfig{image9.pdf}}
    \end{column}
    \begin{column}{0.40\textwidth}
      \begin{itemize}
        \setlength{\itemsep}{0.4em}
        \item AC power transmission principles shown with a simple AC network.
        \item The grid is connected to the load with impedance $Z_L$ via an
              AC cable of length $L_{cable}$.
      \end{itemize}
    \end{column}
  \end{columns}
\end{frame}

\begin{frame}{Transmission line as 2-port units}
  \begin{columns}[T,onlytextwidth]
    \begin{column}{0.54\textwidth}
      \renewcommand{\arraystretch}{1.25}
      \begin{tabular}{@{}l l c c@{}}
        \textbf{Param.} & & \textbf{OHL} & \textbf{Cable (XLPE)} \\
        \hline
        $C'_L$ & nF/km & 10 & 180--420 \\
        $R'_L$ & $\Omega$/km & 0.1--1 & 0.1--1 \\
        $X'_L$ & mH/km & 1 & 0.31--1.42 \\
        $G'_L$ & \textmu S/km & 0.05 & 1 \\
        \hline
      \end{tabular}
      \vspace{0.5em}
      \begin{itemize}
        \setlength{\itemsep}{0.2em}
        \item AC line represented by a $\pi$-section model.
        \item Cables are dominated by capacitance.
        \item OHL are dominated by inductance.
      \end{itemize}
    \end{column}
    \begin{column}{0.42\textwidth}\centering
      \topoimg[width=\linewidth,height=0.4\textheight,keepaspectratio]{\acfig{image10.pdf}}\\[0.4em]
      {\color{atGray}\footnotesize Example cable parameters:\\
      $C'_L=190$\,nF/km, $R'_L=22$\,m$\Omega$/km, $X'_L=1.35$\,mH/km}
    \end{column}
  \end{columns}
\end{frame}

% ---- three load-case slides ---------------------------------
\begin{frame}{Power transmission --- no load connected}
  \begin{columns}[c,onlytextwidth]
    \begin{column}{0.40\textwidth}
      \renewcommand{\arraystretch}{1.2}
      \begin{tabular}{@{}l r@{}}
        \textbf{Param.} & \textbf{Value} \\
        \hline
        $U_{S,LL}$ & 230\,kV (1.0\,p.u.) \\
        $I_S$ & 1060\,A \\
        $P_S$ & 2.2\,MW \\
        $Q_S$ & $-421$\,MVAr \\
        $U_{R,LL}$ & 1.22\,p.u. \\
        $I_R$ & 0\,A \\
        $L_{cable}$ & 60\,mi \\
        \hline
      \end{tabular}
    \end{column}
    \begin{column}{0.56\textwidth}\centering
      \topoimg[width=\linewidth,height=0.5\textheight,keepaspectratio]{\acfig{image11.pdf}}\\[0.3em]
      {\footnotesize Connecting an AC cable creates significant current flow and reactive-power demand.}
    \end{column}
  \end{columns}
\end{frame}

\begin{frame}{Power transmission --- reactive compensation}
  \begin{columns}[c,onlytextwidth]
    \begin{column}{0.40\textwidth}
      \renewcommand{\arraystretch}{1.2}
      \begin{tabular}{@{}l r@{}}
        \textbf{Param.} & \textbf{Value} \\
        \hline
        $U_{S,LL}$ & 230\,kV (1.0\,p.u.) \\
        $I_S$ & 40\,A \\
        $P_S$ & 2.2\,MW \\
        $Q_S$ & $-14.4$\,MVAr \\
        $I_R$ & 830\,A \\
        $Q_R$ & $-277$\,MVAr \\
        $Q_{comp}$ & 400\,MVAr \\
        \hline
      \end{tabular}
    \end{column}
    \begin{column}{0.56\textwidth}\centering
      \topoimg[width=\linewidth,height=0.5\textheight,keepaspectratio]{\acfig{image12.pdf}}\\[0.3em]
      {\footnotesize An uncompensated cable connection is usually not allowed; a shunt reactor $Q_{comp}$ is required.}
    \end{column}
  \end{columns}
\end{frame}

\begin{frame}{Power transmission --- with load}
  \begin{columns}[c,onlytextwidth]
    \begin{column}{0.40\textwidth}
      \renewcommand{\arraystretch}{1.2}
      \begin{tabular}{@{}l r@{}}
        \textbf{Param.} & \textbf{Value} \\
        \hline
        $U_{S,LL}$ & 230\,kV (1.0\,p.u.) \\
        $I_S$ & 1340\,A \\
        $P_S$ & 511\,MW \\
        $Q_S$ & $-150$\,MVAr \\
        $U_{R,LL}$ & 1.06\,p.u. \\
        $I_R$ & 1180\,A \\
        $P_R$ & 500\,MW \\
        \hline
      \end{tabular}
    \end{column}
    \begin{column}{0.56\textwidth}\centering
      \topoimg[width=\linewidth,height=0.5\textheight,keepaspectratio]{\acfig{image13.pdf}}\\[0.3em]
      {\footnotesize Power transmission reduces the voltage rise and the reactive-power demand.}
    \end{column}
  \end{columns}
\end{frame}

\begin{frame}{Power transmission --- reactive power \& current vs.\ active power}
  \begin{columns}[c,onlytextwidth]
    \begin{column}{0.48\textwidth}\centering
      \topoimg[width=\linewidth,height=0.6\textheight,keepaspectratio]{\acfig{image14.png}}\\[0.3em]
      {\footnotesize Reactive power vs.\ active power}
    \end{column}
    \begin{column}{0.48\textwidth}\centering
      \topoimg[width=\linewidth,height=0.6\textheight,keepaspectratio]{\acfig{image16.png}}\\[0.3em]
      {\footnotesize AC current vs.\ active power}
    \end{column}
  \end{columns}
\end{frame}
